In Anbetracht der kurzen Projektdauer ist dennoch ein funktionsfähiger Prototyp entstanden.
Dass dieser noch nicht alle Funktionen vollumfänglich und zur vollsten Zufriedenheit erfüllen wird, war für mich zu erwarten.
Insbesondere bin ich mit meiner Lernkurve innerhalb des Projektes zufrieden. Obwohl ich bisher keine Berührungspunkte im Bereich der mobilen Robotik gehabt habe, wurde es dennoch zusammen mit dem Projektteam erreicht, einen selbstständig navigierenden mobilen Roboter von Grund auf zu entwickeln, der Getränkeuntersetzer spenden kann.

In zukünftigen Projekten würde ich jedoch noch mehr die Schnittstellen zwischen Gewerken in den Blick nehmen und diese frühzeitiger definieren. Das hätte vermutlich einen kompakteren Prototyp ermöglicht und die Übertragung der Software von der Simulation auf die Hardware erleichtert.

Mein Schwerpunkt innerhalb der Umsetzung des Projektes lag in der Gestaltung der mechanischen Komponenten. Hierbei würde ich in Zukunft die Kontaktpunkte mit dem Boden
besser fixieren. Die austauschbaren Seitenwände haben die Montage der Sensoren erleichtert, bei den Motoren jedoch dazu geführt, dass die Bodenplatte nach unten rutscht,
sobald der Deckel des Gehäuses abgenommen wird. Weiterhin würde ich mir mehr Zeit für Tests der Ausgabemechanik einräumen, da erst zum Projektende auffiel, dass 
die vorgesehenen Führungen ein Verkanten der Ausgabemechaniken begünstigen und der Ausschieber zu schmal geraten ist. 

Da trotz allem ein funktionsfähiger Prototyp entstanden ist,
bewerte ich diese Erkenntnisse als lehrreiche, aber nicht erfolgsentscheidend für das Projekt. Daher ist für mich das Projektergebnis insgesamt zufriedenstellend.